\documentclass[10pt]{beamer}

% Tema Limpo e Institucional IFCE Taua
\usetheme{Madrid}
\usecolortheme{whale}

% Pacotes Essenciais
\usepackage[utf8]{inputenc}
\usepackage[T1]{fontenc}
\usepackage[brazilian]{babel}
\usepackage{graphicx}
\usepackage{booktabs}
\usepackage{amsmath}
\usepackage{amssymb}
\usepackage{listings}
\usepackage{xcolor}
\usepackage{tikz}
\usetikzlibrary{shapes.geometric,arrows.meta,positioning,calc}

% Configuracao de Cores e Codigo (SQL e Python)
\definecolor{codegreen}{rgb}{0,0.6,0}
\definecolor{codegray}{rgb}{0.5,0.5,0.5}
\definecolor{codepurple}{rgb}{0.58,0,0.82}
\definecolor{backcolour}{rgb}{0.96,0.97,0.98}

\lstset{
    backgroundcolor=\color{backcolour},
    basicstyle=\ttfamily\scriptsize,
    keywordstyle=\color{blue}\bfseries,
    stringstyle=\color{red!80!black},
    commentstyle=\color{green!60!black},
    numberstyle=\tiny\color{codegray},
    breaklines=true,
    showstringspaces=false,
    tabsize=4,
    frame=single,
    rulecolor=\color{blue!20}
}

% Metadados da Apresentacao
\title[Engenharia de Software -- Aula 25]{Engenharia de Software -- Aula 25}
\subtitle{Sprint 04: Persistência Relacional com SQLite e Padrão Repository}
\author[Prof. Me. Reginaldo Fernandes]{Prof. Me. Reginaldo Pereira Fernandes}
\institute[IFCE Campus Tauá]{
    Instituto Federal de Educação, Ciência e Tecnologia do Ceará -- Campus Tauá\\
    Técnico em Informática para Internet (TSII.210)
}
\date{Semana 1 -- Encontro 2 (100 min)}

\begin{document}

% Slide 1: Titulo
\begin{frame}
    \titlepage
\end{frame}

% Slide 2: Agenda do Encontro
\begin{frame}{Agenda do Encontro 2 (100 Minutos)}
    \tableofcontents
\end{frame}

\section{Mapeamento Objeto-Relacional e DER}

\begin{frame}{Do Modelo de Classes à Persistência em Banco}
    \begin{itemize}
        \item \textbf{Status da Sprint 04:}
        \begin{itemize}
            \item Na Aula 24, modelamos os Casos de Uso e Classes e implementamos \texttt{exceptions.py} e a classe \texttt{Cupom}.
        \end{itemize}
        \item \textbf{O Problema da Memória Volátil:}
        \begin{itemize}
            \item Objetos em memória RAM desaparecem quando o processo encerra.
            \item Os cupons da livraria universitária precisam persistir em disco no banco \textbf{SQLite}!
        \end{itemize}
        \item \textbf{Mapeamento Objeto-Relacional (ORM):}
        \begin{itemize}
            \item Como transformar instâncias de \texttt{Cupom} em registros relacionais?
            \item Como proteger as regras de negócio em duas camadas: código Python E restrições SQL?
        \end{itemize}
    \end{itemize}
\end{frame}

\begin{frame}{Diagrama Entidade-Relacionamento: Esquema Relacional (DER)}
\centering
\resizebox{0.92\textwidth}{!}{%
\begin{tikzpicture}[>=Latex, font=\sffamily]
    % Tabela CUPONS
    \node[draw=blue!80!black, fill=blue!5, thick, rounded corners=2pt, inner sep=3pt] (tbl_cupons) at (-0.2, 0) {
        \begin{tabular}{ll}
        \multicolumn{2}{c}{\textbf{\large TABELA: cupons}} \\
        \hline
        \scriptsize \textbf{id} & \scriptsize INTEGER [PK, AUTO] \\
        \scriptsize \textbf{codigo} & \scriptsize TEXT [UNIQUE, NOT NULL] \\
        \scriptsize tipo\_desconto & \scriptsize TEXT [CHECK Enum] \\
        \scriptsize valor\_desconto & \scriptsize REAL [CHECK > 0] \\
        \scriptsize valor\_minimo & \scriptsize REAL [DEFAULT 0.0] \\
        \scriptsize limite\_usos & \scriptsize INTEGER [CHECK > 0] \\
        \scriptsize usos\_realizados & \scriptsize INTEGER [DEFAULT 0] \\
        \scriptsize validade\_ate & \scriptsize TEXT [ISO8601] \\
        \scriptsize ativo & \scriptsize INTEGER [CHECK 0, 1] \\
        \end{tabular}
    };

    % Tabela PEDIDOS
    \node[draw=green!60!black, fill=green!5, thick, rounded corners=2pt, inner sep=3pt] (tbl_pedidos) at (7.8, 0) {
        \begin{tabular}{ll}
        \multicolumn{2}{c}{\textbf{\large TABELA: pedidos}} \\
        \hline
        \scriptsize \textbf{id} & \scriptsize INTEGER [PK, AUTO] \\
        \scriptsize cliente\_id & \scriptsize INTEGER [FK $\rightarrow$ clientes] \\
        \scriptsize \textbf{cupom\_id} & \scriptsize INTEGER [FK $\rightarrow$ cupons, NULL] \\
        \scriptsize subtotal & \scriptsize REAL [NOT NULL] \\
        \scriptsize desconto & \scriptsize REAL [DEFAULT 0.0] \\
        \scriptsize total & \scriptsize REAL [NOT NULL] \\
        \scriptsize status & \scriptsize TEXT [NOT NULL] \\
        \end{tabular}
    };

    % Conexao de Chave Estrangeira
    \draw[->, very thick, draw=blue!80!black] (tbl_pedidos.west) -- node[midway, above, font=\tiny\bfseries] {FK (cupom\_id)} (tbl_cupons.east);
\end{tikzpicture}
}
\vspace{0.2cm}
\footnotesize
\textit{Associação 0..1: Um pedido pode ou não conter um cupom aplicado.}
\end{frame}

\begin{frame}{Restrições de Integridade (Constraints) no SQLite}
    \begin{block}{Defesa em Profundidade: Por que validar no Banco de Dados?}
    Mesmo que um desenvolvedor esqueça de validar no Python, o motor do SQLite rejeitará dados corrompidos!
    \end{block}
    \vspace{0.2cm}
    \begin{itemize}
        \item \texttt{UNIQUE(codigo)}: Impede duplicidade de nomes de cupons na livraria.
        \item \texttt{CHECK(tipo\_desconto IN ('PERCENTUAL', 'VALOR\_FIXO'))}: Garante consistência com o Enum em Python.
        \item \texttt{CHECK(valor\_desconto > 0)}: Proíbe descontos nulos ou negativos.
        \item \texttt{FOREIGN KEY(cupom\_id) REFERENCES cupons(id)}: Integridade referencial garantida.
    \end{itemize}
\end{frame}

\section{Arquitetura em Camadas e Repositório}

\begin{frame}{Diagrama de Arquitetura: Padrão Repository}
\centering
\resizebox{0.92\textwidth}{!}{%
\begin{tikzpicture}[>=Latex, font=\sffamily]
    % Camada de Dominio
    \node[draw=purple!80!black, fill=purple!8, thick, rounded corners=4pt, text width=3.2cm, align=center] (dom) at (0, 0) {
        \textbf{Camada de Domínio}\\[0.1cm]
        \scriptsize Entidade \texttt{Cupom}\\
        \scriptsize Regras de Negócio Puras\\
        \scriptsize (Sem SQL!)
    };

    % Camada de Repositorio
    \node[draw=blue!80!black, fill=blue!8, thick, rounded corners=4pt, text width=3.4cm, align=center] (repo) at (4.5, 0) {
        \textbf{Camada de Acesso}\\[0.1cm]
        \scriptsize \texttt{CupomRepository}\\
        \scriptsize \texttt{salvar()}, \texttt{buscar()}\\
        \scriptsize \texttt{registrar\_uso()}
    };

    % Camada do Banco de Dados
    \node[draw=green!60!black, fill=green!8, thick, rounded corners=4pt, text width=3.2cm, align=center] (db) at (9.0, 0) {
        \textbf{Banco de Dados}\\[0.1cm]
        \scriptsize Arquivo \texttt{bibliotaua.db}\\
        \scriptsize Tabelas Relacionais\\
        \scriptsize Motor SQLite3
    };

    % Flechas
    \draw[<->, thick, draw=black!70] (dom) -- node[midway, above, font=\tiny\bfseries] {Objetos} (repo);
    \draw[<->, thick, draw=black!70] (repo) -- node[midway, above, font=\tiny\bfseries] {SQL / Tuplas} (db);
\end{tikzpicture}
}
\vspace{0.4cm}
\begin{alertblock}{Princípio de Responsabilidade Única (SRP)}
A entidade \texttt{Cupom} não sabe se os dados vieram de um banco SQLite, PostgreSQL ou arquivo JSON. Toda a persistência é encapsulada pelo repositório!
\end{alertblock}
\end{frame}

\begin{frame}[fragile]{Blindagem contra SQL Injection (Regra Inegociável)}
    \begin{columns}[t]
        \begin{column}{0.48\textwidth}
            \alert{\textbf{Vulnerabilidade Crítica}}\par
            \footnotesize Concatenação de Strings:
\begin{lstlisting}[language=Python]
# NUNCA FACA ISSO!
query = f"SELECT * FROM cupons WHERE codigo = '{cod}'"
cursor.execute(query)
# Permite: ' OR 1=1; DROP TABLE...
\end{lstlisting}
        \end{column}
        \begin{column}{0.48\textwidth}
            \textbf{Padrão Seguro (Parametrizado)}\par
            \footnotesize Parâmetros Posicionais:
\begin{lstlisting}[language=Python]
# SEGURO: Parametros com ?
query = "SELECT * FROM cupons WHERE codigo = ?"
cursor.execute(query, (cod,))
# O driver trata caracteres especiais como dado puro!
\end{lstlisting}
        \end{column}
    \end{columns}
\end{frame}

\section{Implementação: DDL SQLite}

\begin{frame}[fragile]{Mão na Massa: DDL da Tabela \texttt{cupons} (\texttt{database.py})}
\begin{lstlisting}[language=SQL]
CREATE TABLE IF NOT EXISTS cupons (
    id INTEGER PRIMARY KEY AUTOINCREMENT,
    codigo TEXT NOT NULL UNIQUE,
    tipo_desconto TEXT NOT NULL CHECK(tipo_desconto IN ('PERCENTUAL', 'VALOR_FIXO')),
    valor_desconto REAL NOT NULL CHECK(valor_desconto > 0),
    valor_minimo_pedido REAL NOT NULL DEFAULT 0.0 CHECK(valor_minimo_pedido >= 0),
    limite_usos INTEGER NOT NULL DEFAULT 100 CHECK(limite_usos > 0),
    usos_realizados INTEGER NOT NULL DEFAULT 0 CHECK(usos_realizados >= 0),
    validade_ate TEXT NOT NULL,
    ativo INTEGER NOT NULL DEFAULT 1 CHECK(ativo IN (0, 1))
);

CREATE INDEX IF NOT EXISTS idx_cupons_codigo ON cupons(codigo);
\end{lstlisting}
\end{frame}

\section{Implementação: Repositório CupomRepository}

\begin{frame}[fragile]{Mão na Massa: \texttt{CupomRepository} (Persistência e Busca)}
\begin{lstlisting}[language=Python]
class CupomRepository:
    def __init__(self, conn: sqlite3.Connection):
        self.conn = conn

    def salvar(self, cupom: Cupom) -> Cupom:
        cursor = self.conn.cursor()
        cursor.execute("""
            INSERT INTO cupons (codigo, tipo_desconto, valor_desconto, valor_minimo_pedido, limite_usos, validade_ate, ativo)
            VALUES (?, ?, ?, ?, ?, ?, ?)
        """, (cupom.codigo, cupom.tipo_desconto.value, cupom.valor_desconto,
              cupom.valor_minimo_pedido, cupom.limite_usos, cupom.validade_ate.isoformat(), 1 if cupom.ativo else 0))
        self.conn.commit()
        cupom.id = cursor.lastrowid
        return cupom

    def buscar_por_codigo(self, codigo: str) -> Optional[Cupom]:
        cursor = self.conn.cursor()
        cursor.execute("SELECT * FROM cupons WHERE codigo = ?", (codigo.strip().upper(),))
        row = cursor.fetchone()
        return self._row_para_objeto(row) if row else None
\end{lstlisting}
\end{frame}

\begin{frame}[fragile]{Mão na Massa: Incremento Atômico Seguro (\texttt{registrar\_uso})}
\begin{lstlisting}[language=Python]
    def registrar_uso(self, cupom_id: int) -> bool:
        """Incrementa atomicamente os usos realizados sem risco de corrida."""
        cursor = self.conn.cursor()
        cursor.execute("""
            UPDATE cupons 
            SET usos_realizados = usos_realizados + 1
            WHERE id = ? AND usos_realizados < limite_usos AND ativo = 1
        """, (cupom_id,))
        self.conn.commit()
        if cursor.rowcount == 0:
            raise CupomEsgotadoError(0)
        return True
\end{lstlisting}
    \begin{alertblock}{Atomicidade e Proteção contra Concorrência}
    A cláusula \texttt{WHERE usos\_realizados < limite\_usos} impede que dois clientes utilizem o último cupom simultaneamente!
    \end{alertblock}
\end{frame}

\section{Rito Ágil e Próximos Passos}

\begin{frame}{Rito Ágil: Registro da Daily 2}
    \begin{block}{Anotação para o Relatório da Sprint no \texttt{guia\_desenvolvedor.md}}
    \begin{itemize}
        \item \textbf{O que realizei hoje?}
        \begin{itemize}
            \item Projetei o Diagrama Entidade-Relacionamento (DER) integrando \texttt{cupons} e \texttt{pedidos}.
            \item Escrevi o script DDL com restrições \texttt{CHECK}, \texttt{UNIQUE} e índices em \texttt{database.py}.
            \item Implementei a classe \texttt{CupomRepository} com queries parametrizadas blindadas.
            \item Implementei o incremento atômico de utilizações \texttt{registrar\_uso}.
        \end{itemize}
        \item \textbf{O que farei no Encontro 3 (Aula 26)?}
        \begin{itemize}
            \item Analisar o Diagrama de Sequência do checkout.
            \item Implementar o serviço de cálculo \texttt{CalculadoraDescontoService}.
            \item Integrar o cupom ao agregado \texttt{Pedido} e implementar o \texttt{CheckoutService}.
        \end{itemize}
        \item \textbf{Impedimentos?} Nenhum impedimento.
    \end{itemize}
    \end{block}
\end{frame}

\begin{frame}{Síntese do Encontro 2}
    \begin{itemize}
        \item \textbf{Conquistas de Hoje:}
        \begin{itemize}
            \item Mapeamento ORM concluído com sucesso.
            \item Persistência relacional protegida em SQLite.
            \item Camada de dados isolada da camada de regras de negócio.
        \end{itemize}
        \item \textbf{Para o Encontro 3 (Aula 26):}
        \begin{itemize}
            \item Orquestração do Negócio: \textbf{Diagrama de Sequência UML}.
            \item Cálculo de desconto e finalização de pedidos.
            \item Baixa no estoque de livros e transações completas.
        \end{itemize}
    \end{itemize}
    \vspace{0.4cm}
    \centering
    \textbf{Excelente progresso! A persistência do BiblioTauá está blindada!}
\end{frame}

\end{document}
